% !TEX encoding = UTF-8 Unicode
\documentclass[11pt,a4paper]{article}

\usepackage[utf8]{inputenc}
\usepackage[T1]{fontenc}
\usepackage{geometry}
\usepackage{hyperref}
\usepackage{enumitem}
\usepackage{booktabs}
\usepackage{array}
\usepackage{tabularx}
\usepackage{xcolor}
\usepackage{fancyhdr}
\usepackage{titlesec}
\usepackage{parskip}
\usepackage{listings}
\usepackage{tcolorbox}
\tcbuselibrary{breakable, skins, listings}

\geometry{a4paper, top=2.2cm, bottom=2.2cm, left=2.5cm, right=2.5cm}

\definecolor{accentblue}{RGB}{30, 80, 160}
\definecolor{promptbg}{RGB}{242, 244, 248}
\definecolor{promptborder}{RGB}{30, 80, 160}
\definecolor{warnbg}{RGB}{255, 248, 235}
\definecolor{warnborder}{RGB}{190, 80, 0}
\definecolor{tipbg}{RGB}{240, 248, 240}
\definecolor{tipborder}{RGB}{50, 130, 60}
\definecolor{darkgray}{RGB}{70, 70, 70}
\definecolor{medgray}{RGB}{130, 130, 130}
\definecolor{cheatbg}{RGB}{250, 250, 252}
\definecolor{cheatborder}{RGB}{30, 80, 160}
\definecolor{labelcolor}{RGB}{40, 40, 100}

\pagestyle{fancy}
\fancyhf{}
\fancyhead[L]{\small\textcolor{medgray}{Prompt engineering a gyakorlatban}}
\fancyhead[R]{\small\textcolor{medgray}{Sapientia EMTE -- dr.\ Pál László}}
\fancyfoot[C]{\small\textcolor{medgray}{\thepage}}
\renewcommand{\headrulewidth}{0.3pt}

\titleformat{\section}{\large\bfseries\color{accentblue}}{\thesection}{0.6em}{}%
  [\vspace{-2pt}\textcolor{accentblue}{\rule{\linewidth}{0.6pt}}]
\setcounter{secnumdepth}{0}
\titleformat{\subsection}{\normalsize\bfseries\color{darkgray}}{}{0em}{}
\titlespacing{\section}{0pt}{18pt}{8pt}
\titlespacing{\subsection}{0pt}{12pt}{4pt}

% lstlisting stílus a prompt boxokhoz -- a literate kezeli a magyar karaktereket
\lstdefinestyle{promptlst}{
  basicstyle=\small\ttfamily\color{black},
  breaklines=true,
  breakatwhitespace=false,
  frame=none,
  columns=fullflexible,
  keepspaces=true,
  showstringspaces=false,
  literate=
    {á}{{\'{a}}}1 {é}{{\'{e}}}1 {í}{{\'{i}}}1 {ó}{{\'{o}}}1 {ö}{{\"o}}1
    {ő}{{\H{o}}}1 {ú}{{\'{u}}}1 {ü}{{\"u}}1 {ű}{{\H{u}}}1
    {Á}{{\'{A}}}1 {É}{{\'{E}}}1 {Í}{{\'{I}}}1 {Ó}{{\'{O}}}1 {Ö}{{\"O}}1
    {Ő}{{\H{O}}}1 {Ú}{{\'{U}}}1 {Ü}{{\"U}}1 {Ű}{{\H{U}}}1,
}

\tcbset{
  promptstyle/.style={
    breakable, enhanced,
    colback=promptbg, colframe=promptborder,
    boxrule=1pt, arc=3pt,
    left=10pt, right=10pt, top=10pt, bottom=8pt,
    attach boxed title to top left={xshift=10pt, yshift=-\tcboxedtitleheight/2},
    boxed title style={
      colback=promptborder, colframe=promptborder,
      arc=2pt, boxrule=0pt,
      left=6pt, right=6pt, top=2pt, bottom=2pt
    },
    coltitle=white, fonttitle=\footnotesize\bfseries\sffamily,
  },
  warnstyle/.style={
    breakable, colback=warnbg, colframe=warnborder,
    boxrule=0.5pt, leftrule=3pt, arc=2pt,
    left=8pt, right=8pt, top=5pt, bottom=5pt,
    fontupper=\small\sffamily,
  },
  tipstyle/.style={
    breakable, colback=tipbg, colframe=tipborder,
    boxrule=0.5pt, leftrule=3pt, arc=2pt,
    left=8pt, right=8pt, top=5pt, bottom=5pt,
    fontupper=\small\sffamily,
  },
  cheatstyle/.style={
    breakable, enhanced, colback=cheatbg, colframe=cheatborder,
    boxrule=0.8pt, arc=3pt,
    left=10pt, right=10pt, top=9pt, bottom=7pt,
    attach boxed title to top left={xshift=10pt, yshift=-\tcboxedtitleheight/2},
    boxed title style={
      colback=cheatborder, colframe=cheatborder,
      arc=2pt, boxrule=0pt,
      left=6pt, right=6pt, top=2pt, bottom=2pt
    },
    coltitle=white, fonttitle=\footnotesize\bfseries\sffamily,
  },
}

% Prompt box: newtcblisting a listings integrációhoz
\newtcblisting{prompt}[1][Prompt]{
  promptstyle, title=#1,
  listing only,
  listing options={style=promptlst},
}

\newcommand{\warn}[1]{%
  \vspace{3pt}%
  \begin{tcolorbox}[warnstyle]%
  \textbf{\textcolor{warnborder}{Figyelem:}} #1%
  \end{tcolorbox}\vspace{2pt}}

\newcommand{\tip}[1]{%
  \vspace{3pt}%
  \begin{tcolorbox}[tipstyle]%
  \textbf{\textcolor{tipborder}{Mikor hasznos?}} #1%
  \end{tcolorbox}\vspace{2pt}}

\newcommand{\lbl}[1]{\vspace{6pt}\noindent{\small\bfseries\textcolor{labelcolor}{#1}}\par\smallskip}

\hypersetup{colorlinks=true, linkcolor=accentblue, urlcolor=accentblue}

\begin{document}

\thispagestyle{empty}
\begin{center}
  {\LARGE\bfseries Prompt engineering a gyakorlatban}\\[5pt]
  {\normalsize Generatív AI és üzleti alkalmazások modul}\\[1pt]
  {\small\color{medgray} dr.\ Pál László, Sapientia EMTE, 2026. március 6--7.}
\end{center}

\vspace{10pt}
\noindent
\begin{tabularx}{\linewidth}{lXr}
\toprule
\textbf{Blokk} & \textbf{Tartalom} & \textbf{Idő} \\
\midrule
1.\ rész & RTFC módszertan (6 lépés + delimiterek) & $\sim$25 perc \\
2.\ rész & 5 haladó technika (Few-shot, CoT, Persona, Negatív, Reflexion) & $\sim$25 perc \\
3.\ rész & Zárófeladat (dokumentum-feldolgozás) & $\sim$10 perc \\
\midrule
\multicolumn{2}{l}{\textbf{Teljes blokk}} & \textbf{$\sim$60 perc} \\
\bottomrule
\end{tabularx}

\smallskip
\noindent\textbf{Eszköz:} Gemini (\texttt{gemini.google.com})

\vspace{8pt}

%=============================================================
\section{1. RÉSZ -- RTFC módszertan}
%=============================================================

\noindent\textit{Cél: Megérteni, hogyan javítja a szerep, a feladat, a formátum és a kontextus pontos meghatározása a prompt minőségét.}

\bigskip
\noindent Az \textbf{RTFC} keretrendszer:
\smallskip

\begin{tabularx}{\linewidth}{@{}llX@{}}
\toprule
\textbf{Elem} & \textbf{Fogalom} & \textbf{Jelentés} \\
\midrule
\textbf{R} & Role -- Szerep    & Milyen szerepkörben válaszoljon az AI? \\
\textbf{T} & Task -- Feladat   & Mit kell csinálnia? \\
\textbf{F} & Format -- Formátum & Milyen formában, milyen terjedelemben? \\
\textbf{C} & Context -- Kontextus & Milyen háttérinformáció áll rendelkezésre? \\
\bottomrule
\end{tabularx}

\bigskip

%--- 1. lépés ---
\subsection{1. lépés -- Egyszerű prompt}

\noindent Ha nem mondunk semmit, az AI mindent maga dönt el -- témát, hangnemet, hosszt. Próbáljuk ki.

\begin{prompt}
Írj egy emailt egy ügyfélnek.
\end{prompt}

\noindent Mindenki más eredményt kap. Ez a prompt engineering kiindulópontja.

\warn{Hallucináció: az AI magabiztosan közölhet nem létező vagy pontatlan tényeket és adatokat. A konkrét számokat és hivatkozásokat minden esetben ellenőrizzük.}

%--- 2. lépés ---
\subsection{2. lépés -- Kontextus hozzáadása}

\noindent Ugyanaz a kérés, de most megadjuk a kontextust. Vizsgáljuk meg, hogyan változik az eredmény.

\begin{prompt}
Írj emailt egy ügyfélnek, aki késve kapta meg a csomagját.
\end{prompt}

\noindent A szituáció pontosítása javítja az eredményt, ugyanakkor a hangnemet, a terjedelmet és a konkrét ajánlást továbbra is a modell alakítja.

%--- 3. lépés ---
\subsection{3. lépés -- Szerep megadása (Role)}

\noindent Ugyanaz a feladat két különböző szerepkörben megfogalmazva. A Role önmagában is jelentősen befolyásolja a hangnemet.

\medskip\noindent\textbf{1. Szerepkör:} Barátságos ügyfélszolgálatos munkatárs

\begin{prompt}[Prompt -- 1. szerepkör]
Te egy barátságos ügyfélszolgálatos munkatárs vagy.
Írj emailt egy ügyfélnek, aki késve kapta meg a csomagját.
\end{prompt}

\noindent\textbf{2. Szerepkör:} Szigorú, jogi szemléletű ügyfélszolgálati vezető

\begin{prompt}[Prompt -- 2. szerepkör]
Te egy szigorú, jogi szemléletű ügyfélszolgálati vezető vagy.
Írj emailt egy ügyfélnek, aki késve kapta meg a csomagját.
\end{prompt}

\noindent Azonos feladat, teljesen más kimenet.

%--- 4. lépés ---
\subsection{4. lépés -- Formátum és határolók (delimiters)}

\noindent A határolók (\texttt{\#\#\#}, \texttt{"""}) segítenek az AI-nak elkülöníteni az utasítást a tartalomtól. Különösen hasznos, ha hosszabb szöveget, adatot illesztünk be a promptba.

\lbl{Határolók nélkül}
\begin{prompt}
Te egy barátságos ügyfélszolgálatos munkatárs vagy. Írj emailt
egy ügyfélnek, aki késve kapta meg a csomagját. Max. 5 mondat,
legyen bocsánatkérés és ajánlj fel 10% kedvezményt.
\end{prompt}

\lbl{Határolókkal strukturálva}
\begin{prompt}[Prompt -- határolókkal]
Te egy barátságos ügyfélszolgálatos munkatárs vagy.

### Feladat
Írj emailt egy ügyfélnek, aki késve kapta meg a csomagját.

### Kontextus
"""
Csomag száma: RO987654
Várható: 2026.02.15.   
Tényleges: 2026.03.02.
Vevő: Nagy István, Kolozsvár
"""

### Formátum
- maximum 5 mondat
- bocsánatkérés kötelező
- ajánlj fel 10% kedvezményt a következő rendelésre
- zárd kérdéssel
\end{prompt}

\tip{A \texttt{"""} jelöli az adatblokk határait -- hasznos, ha saját anyagot (önéletrajz, cégleírás, termékadatok) illesztünk be a promptba.}

%--- 5. lépés ---
\subsection{5. lépés -- Iteráció}

\noindent A Gemini emlékszik az előző üzenetre — nem kell újraírni az egész promptot. Küldjünk egy új üzenetet az előző válasz finomítására.

\begin{prompt}[Finomítás]
Legyen formálisabb.
\end{prompt}
\begin{prompt}[Finomítás]
Rövidítsd 3 mondatra.
\end{prompt}
\begin{prompt}[Finomítás]
Add hozzá, hogy 48 órán belül visszajelzünk a kompenzációról.
\end{prompt}



%--- 6. lépés ---
\subsection{6. lépés -- Önálló feladat}

\noindent Alkalmazzuk az eddig tanultakat önállóan: RTFC keretrendszer és határolók, új témán. Mindenki választ egy változatot.

\begin{prompt}[Önálló feladat -- A változat]
Termék: vezeték nélküli fülhallgató
Cél: rövid termékleírás webshop stílusban
\end{prompt}

\begin{prompt}[Önálló feladat -- B változat]
Szituáció: állásinterjúra készülsz egy csíkszeredai marketingcégnél
Cél: rövid bemutatkozó email a HR-esnek
\end{prompt}

\noindent Tippek: adj meg szerepkört (R), fogalmazd meg pontosan a feladatot (T), határozd meg a formátumot -- max.\ hány mondat, milyen hangnem (F), és illeszd be a releváns háttéradatokat határolókkal (C).

%=============================================================
\section{2. RÉSZ -- Haladó technikák}
%=============================================================

\noindent\textit{Cél: Az RTFC-n túl négy technika, amelyek összetettebb feladatokhoz szükségesek.}

\bigskip

%--- Few-shot ---
\subsection{A) Few-shot prompting -- tanítás példákkal}

Ha megmutatjuk az AI-nak a kívánt stílust, abból képes következtetni a válasz formájára.

\begin{prompt}
Íme két rövid kommentár, tárgyilagos piaci elemzői stílusban:

1. "A lakáspiac stagnált."
   --> "A tranzakciók volumene nem változott, a kereslet továbbra is óvatos."

2. "Az infláció gyorsult."
   --> "A drágulás fő oka az energiaárak emelkedése."

Írj hasonló stílusú kommentárt erre:
"A kiskereskedelmi forgalom 2026 januárban 2.1%-kal csökkent
az előző év azonos hónapjához képest."

Alkalmazd az RTFC keretrendszert:
R = piaci elemző
T = tárgyilagos kommentár
F = max. 2 mondat
C = romániai kiskereskedelmi adat
\end{prompt}

\lbl{Alternatív feladat -- szarkasztikus tech poszt}

\begin{prompt}[Prompt -- alternatív]
Íme két szarkasztikus tech-hír poszt:

Hír: Az új okosszemüveg akkumulátora csak 2 órát bír.
Poszt: A jövő itt van -- vehetünk egy méregdrága szemüveget,
       amivel pont megnézhetünk egy rövid videót.

Hír: Egy AI megnyert egy művészeti versenyt a művészek tudta nélkül.
Poszt: A robotok már nemcsak a munkánkat, hanem a kreativitásunkat
       is próbára teszik. Gratulálunk a győztesnek.

Írj hasonló stílusú posztot:
Hír: A legújabb szoftverfrissítés véletlenül törölte a felhasználók
     mentett jelszavait.
\end{prompt}

\tip{Ha a stílus fontosabb, mint a tartalom -- könnyebb megmutatni, mint szavakkal körülírni.}

%--- CoT ---
\subsection{B) Chain of Thought (CoT) -- lépésekben gondolkodj}

Strukturált lépésekre bontva jobb elemzést kapunk, mint nyílt kérdéssel.

\begin{prompt}[Prompt -- CoT nélkül]
Egy erdélyi ruházati webshop havi eladási adatai:
Január: 180 db   Február: 210 db   Március: 165 db
Átlagos kosárérték: 380 lej

Elemezd az adatokat és adj javaslatot áprilisra.
\end{prompt}

\begin{prompt}[Prompt -- CoT használatával]
Egy erdélyi ruházati webshop havi eladási adatai:
Január: 180 db   Február: 210 db   Március: 165 db
Átlagos kosárérték: 380 lej

Gondolkodj lépésről lépésre:
1. Melyik hónapban volt a legnagyobb / legkisebb bevétel?
2. Mi lehet a márciusi visszaesés oka (szezon, akció, versenytárs)?
3. Milyen marketing lépést javasolsz áprilisra?

Minden lépés max. 2 mondat.
\end{prompt}

\lbl{Alternatív feladat -- részvényelemzés}

\begin{prompt}[Prompt -- alternatív]
Három technológiai részvény teljesítménye 2025. december 31. --
2026. február 27. között (záróárak alapján, Yahoo Finance):
- Apple (AAPL):  271.86 USD --> 264.18 USD = -2.8%
- Tesla (TSLA):  449.72 USD --> 402.51 USD = -10.5%
- NVIDIA (NVDA): 186.50 USD --> 177.19 USD = -5.0%
(forrás: Yahoo Finance, 2026. március 2.)

Gondolkodj lépésről lépésre:
1. Melyik részvény teljesített a legjobban / legrosszabbul?
2. Mi lehetett a fő oka a Tesla ~10.5%-os csökkenésének 2026 elején?
3. Ha 10 000 lejt egyenlő arányban fektettél volna be mindhárom
   részvénybe 2025. december 31-én (USD/LEJ árfolyamot vedd
   figyelembe), mennyi lenne a portfólió értéke február 27-én?

Minden lépés max. 2 mondat. Táblázattal ahol lehet.

Tipp -- alkalmazd az RTFC keretrendszert:
R = befektető
T = rövid elemzés + allokáció
F = max. 2 mondat/lépés
C = technológiai részvények
\end{prompt}

\tip{Döntéstámogatásnál és elemzésnél, ahol az indoklás legalább annyit számít, mint a végeredmény.}

%--- Persona ---
\subsection{C) Szerepkör beállítása (Persona)}

Az AI végig egy adott \textit{personaként} viselkedik – különböző nézőpontok megjelenítésére alkalmas.

\medskip
\noindent\textbf{1. Persona:} Konzervatív banki hitelbíráló\\
-- tények, kockázatok.

\noindent\textbf{2. Persona:} Lelkes EdTech befektető\\
-- növekedési potenciál.

\begin{prompt}
Üzleti ötlet: Online használt tankönyv piactér egyetemistáknak  
(kezdő befektetés: 25 000 lej marketingre + weboldalra)

Értékeld az ötletet kétféle szerepkörből:

1. Szerepkör: konzervatív banki hitelbíráló (hagyományos banki szemlélet)
   -- tények, kockázatok.

2. Szerepkör: lelkes EdTech befektető (startup / oktatástechnológia)
   -- növekedési potenciál.
\end{prompt}

\tip{Több nézőpont egyidejű megszerzéséhez, döntések előkészítéséhez.}

%--- Negatív ---
\subsection{D) Negatív utasítások}

Megmondjuk, mit \emph{ne} csináljon az AI. Sokszor hatékonyabb, mint hosszú pozitív leírás.

\begin{prompt}[Prompt -- korlátok nélkül]
Írj termékleírást egy okostelefonról.
\end{prompt}

\begin{prompt}[Prompt -- negatív utasításokkal]
R = termékleíró
T = okostelefon termékleírás
F = max. 4 mondat, marketing klisék nélkül
C = webshop termékbemutató

Írj termékleírást egy okostelefonról.
NE használd: "forradalmi", "game-changer", "diszruptív",
             "exponenciális növekedés".
NE tegyél bele árat.
NE kezdd azzal, hogy "Bemutatkozik..."
\end{prompt}

\tip{Negatív utasításokkal pontosabban irányíthatjuk a válasz tartalmát és stílusát.}


%--- Reflexion ---
\subsection{E) Reflexion -- ellenőriztesd magával az AI-val}

Az AI válaszai lehetnek hibásak. A Reflexion technikával megkérjük, hogy
válaszadás után ellenőrizze saját megoldását, keressen hibákat vagy
ellentmondásokat, és adjon javított verziót.

\begin{prompt}[Prompt -- Reflexion nélkül]
Egy webshop januárban 200, februárban 180, márciusban 260
terméket adott el. Mekkora volt az átlagos havi növekedési ütem?
\end{prompt}

\noindent Reflexionnal:

\begin{prompt}[Prompt -- Reflexionnal]

R = elemző
T = átlagos havi növekedési ütem számítása és ellenőrzése
F = számítás + önellenőrzés
C = webshop eladási adatok

Egy webshop januárban 200, februárban 180, márciusban 260
terméket adott el. Mekkora volt az átlagos havi növekedési ütem?

### Reflexion
Miután megválaszoltad, nézd át saját magad:
- helyes-e a számítás menete
- vannak-e kerekítési vagy logikai hibák
- a növekedési ütem értelmezése helyes-e

Javítsd ki az esetleges hibákat, és add meg a végleges választ.
\end{prompt}

\tip{A Reflexion különösen hasznos számításoknál és elemzéseknél, ahol a
pontosság kritikus. Az önellenőrzés javítja a válasz minőségét és
megbízhatóságát.}

\section{3. RÉSZ -- Zárófeladat: dokumentumfeldolgozás}

\noindent\textit{Cél: Az összes tanult technika önálló alkalmazása valós
dokumentumokon. Minden hallgató feltölti a választott dokumentumot a Gemini rendszerébe,
majd sorban elvégzi az alábbi feladatokat.}

\bigskip
\noindent\textbf{Válasszon egyet az alábbi három dokumentum közül:}

\begin{itemize}[noitemsep, topsep=4pt]
  \item[\textbf{A}] Nistor \& Zadobrischi: \textit{The Scale-Up of
        E-Commerce in Romania Generated by the Pandemic, Automation,
        and Artificial Intelligence.} Telecom 2024, 5, 680--705.
        \textit{(angol, kutatási cikk)}
  \item[\textbf{B}] Póka \& Lányi: \textit{Az utolsó száz méter
        kihívásai az e-kereskedelem logisztikában.}
        Acta Periodica XXVI, 2022.
        \textit{(magyar, kutatási cikk)}
  \item[\textbf{C}] \textit{Záróvizsgák megszervezésének módszertana
        a Sapientia EMTE-n,} 2025/26.\ tanév.
        \textit{(magyar, intézményi szabályzat)}
\end{itemize}

\bigskip\noindent\rule{\linewidth}{0.4pt}\bigskip

%--------------------------------------------------------------
\subsection{A változat -- Telecom cikk (Nistor \& Zadobrischi)}
%--------------------------------------------------------------

\noindent\textit{Téma: A román e-kereskedelem növekedése, valamint az AI és az automatizálás hatása.
Angol nyelvű, 26 oldalas kutatási cikk, konkrét piaci adatokkal
és matematikai modellel.}

\bigskip

\subsubsection*{1.\ feladat -- Alapösszefoglalás (RTFC)}

\begin{prompt}
Te egy tapasztalt üzleti elemző vagy.
### Feladat
Foglald össze ezt a dokumentumot 5 mondatban.
### Formátum
- Folyó szöveg, felsorolás nélkül.
- Célközönség: nem szakértő olvasó, aki soha nem látta a cikket.
\end{prompt}

\noindent\textbf{Kérdés:} Felismerhetők az RTFC elemei?
Hol található a Role, a Task és a Format -- és mi jelenti a Contextet?
\textit{(A kontextus maga a feltöltött PDF.)}

\bigskip

\subsubsection*{2.\ feladat -- Strukturált kinyerés}

\begin{prompt}
A feltöltött dokumentumból nyerd ki az alábbi információkat
KIZÁRÓLAG Markdown-táblázat formájában (ne írj bevezető szöveget):
- Szerzők és intézmény
- Fő kutatási kérdés
- Alkalmazott módszertan
- A 3 legfontosabb eredmény (konkrét számadatokkal)
- Korlátok / jövőbeli irányok

Ha valamely információ nem szerepel --> "Nem szerepel a dokumentumban"
\end{prompt}

\noindent\textit{Hallucinációteszt:} A táblázatban szereplő számadatokat
ellenőrizzük az eredeti cikkben! Például a 6,5 milliárd eurós piaci értéket,
a 38,4\%-os növekedést vagy a 309,5\%-os szépségipari bővülést.

\bigskip

\subsubsection*{3.\ feladat -- Lépésről lépésre elemzés}

\begin{prompt}
Te egy romániai e-kereskedelmi tanácsadó vagy.
A feltöltött cikk adatai alapján gondolkodj lépésről lépésre:

1. Melyik termékkategória nőtt a legnagyobb mértékben 2020-ban,
   és mi lehetett ennek az oka?
2. Mit jelent az, hogy a román webshopok bevételük akár 50\%-át
   is elveszíthetik az automatizálás hiánya miatt? Reális ez a becslés?
3. Mit javasolnál egy erdélyi kis webshopnak 2026-ra a cikk alapján?

Minden pontra legfeljebb 2 mondatban válaszolj.
\end{prompt}

\noindent\textit{Figyeljük meg:} a „gondolkodj lépésről lépésre”
utasítás nélkül az AI valószínűleg egy összefoglaló jellegű választ adna.
Próbáljuk ki anélkül is -- érzékelhető a különbség?

\bigskip\noindent\rule{\linewidth}{0.4pt}\bigskip

%--------------------------------------------------------------
\subsection{B változat -- Last-mile cikk (Póka \& Lányi)}
%--------------------------------------------------------------

\noindent\textit{Téma: Az e-kereskedelem utolsó mérföldjének logisztikai
kihívásai Magyarországon. Magyar nyelvű, 16 oldalas tudományos cikk,
603 fős empirikus felméréssel.}

\bigskip

\subsubsection*{1.\ feladat -- Alapösszefoglalás (RTFC)}

\begin{prompt}
Te egy tapasztalt üzleti elemző vagy.
### Feladat
Foglald össze ezt a dokumentumot 5 mondatban.
### Formátum
- Folyó szöveg, felsorolás nélkül.
- Célközönség: nem szakértő olvasó, aki soha nem látta a cikket.
\end{prompt}

\noindent\textit{Összehasonlítás:} Ugyanaz a prompt, mint az A
változatnál -- a kimenet mégis más jellegű lesz. Mi ennek az oka?

\bigskip

\subsubsection*{2.\ feladat -- Strukturált kinyerés}

\begin{prompt}
A feltöltött dokumentumból nyerd ki az alábbi információkat
KIZÁRÓLAG Markdown-táblázat formájában (ne írj bevezető szöveget):
- Szerzők és intézmény
- Fő kutatási kérdés
- Alkalmazott módszertan
- A 3 legfontosabb eredmény (konkrét számadatokkal)
- Korlátok / jövőbeli irányok

Ha valamely információ nem szerepel --> "Nem szerepel a dokumentumban"
\end{prompt}

\noindent\textit{Érdekesség:} A cikk tartalmaz egy 603 fős felmérést
a zöld megoldásokért való fizetési hajlandóságról.
Kinyeri-e az AI ezt az adatot -- és pontosan idézi-e?

\bigskip

\subsubsection*{3.\ feladat -- Kettős persona}

\begin{prompt}
Értékeld a feltöltött cikk főbb megállapításait két különböző szerepkörből:

1. Szerepkör: szkeptikus akadémiai lektor
   -- Milyen módszertani gyengeségeket azonosítasz?
      Mennyire tekinthetők megalapozottnak az állítások?

2. Szerepkör: romániai logisztikai startup alapítója
   -- Milyen üzleti lehetőségeket látsz a cikk alapján?
      Min kezdenél el dolgozni már holnap?

Mindkét szerepkörben legfeljebb 3 mondatban válaszolj.
\end{prompt}

\noindent\textit{Figyeljük meg:} Ugyanarról a szövegről két teljesen
különböző értékelés születik attól függően, milyen szerepkört adunk
az AI-nak. Ez jól mutatja a persona technika hatását.

\bigskip\noindent\rule{\linewidth}{0.4pt}\bigskip

%--------------------------------------------------------------
\subsection{C változat -- Záróvizsga-szabályzat (Sapientia EMTE)}
%--------------------------------------------------------------

\noindent\textit{Téma: A Sapientia EMTE záróvizsgáinak megszervezési
módszertana a 2025/26.\ tanévben. Magyar nyelvű, 11 oldalas intézményi
szabályzat, amely minden hallgatót közvetlenül érint.}

\bigskip

\subsubsection*{1.\ feladat -- Alapösszefoglalás (RTFC)}

\begin{prompt}
Te egy tapasztalt üzleti elemző vagy.
### Feladat
Foglald össze ezt a dokumentumot 5 mondatban.
### Formátum
- Folyó szöveg, felsorolás nélkül.
- Célközönség: nem szakértő olvasó, aki soha nem látta a dokumentumot.
\end{prompt}

\noindent\textit{Összehasonlítás:} Ismét ugyanaz a prompt, de most
egy szabályzati dokumentum esetében. Hogyan változik a kimenet stílusa
a kutatási cikkekhez képest?

\bigskip

\subsubsection*{2.\ feladat -- Strukturált kinyerés + negatív utasítások}

\begin{prompt}
Te egy hallgatói tanácsadó vagy.
### Feladat
A feltöltött záróvizsga-szabályzatból nyerd ki az alábbi információkat
KIZÁRÓLAG Markdown-táblázat formájában.
### Tartalom
- Záróvizsgára bocsátás feltételei
- Beiratkozáshoz szükséges dokumentumok
- A jegyszámítás módja (I. és II. próba)
- Fellebbezés határideje és menete
- Sikertelen vizsga esetén a teendők
### Korlátok
NE írj bevezető szöveget a táblázat elé.
NE idézz cikkelyszámokat -- csak a lényeget fogalmazd meg.
Ha valamely információ nem szerepel --> "Nem szerepel a szabályzatban"
\end{prompt}

\noindent\textit{Figyeljük meg a negatív utasításokat:}
a „NE idézz cikkelyszámokat” megkötés nélkül az AI válasza
valószínűleg tele lenne „5.\ cikk (2) bekezdés” típusú hivatkozásokkal.
Próbáljuk ki anélkül is!

\bigskip

\subsubsection*{3.\ feladat -- Few-shot: FAQ stílus}

\begin{prompt}
Íme két példa arra, hogyan kell hallgatói kérdésre tömören válaszolni
a szabályzat alapján:

K: "Mikor kell leadni a szakdolgozatot?"
V: "A témavezető által elfogadott dolgozatot a záróvizsgára való
    beiratkozással egy időben, elektronikus formában kell benyújtani
    a dékáni hivatalba."

K: "Mi történik, ha megbukom a záróvizsgán?"
V: "Egy következő vizsgaidőszakban újra lehet jelentkezni.
    Ha az átmenő részjegy megmarad, azt a bizottság az újabb vizsgán
    elismeri."

Most írj ugyanilyen stílusú válaszokat az alábbi kérdésekre
a feltöltött szabályzat alapján:
1. "Kötelező a nyelvvizsga a záróvizsgához?"
2. "Hogyan számítják ki a záróvizsga végső jegyét?"
3. "Lehet angolul írni a szakdolgozatot?"
\end{prompt}

\noindent\textit{A few-shot technika itt azt mutatja meg az AI számára,
hogy tömör, közvetlen, hallgatóbarát stílusban válaszoljon -- ne
szabályzati nyelvezetben. A két példapár ezt a mintát szemlélteti.}

\bigskip\noindent\rule{\linewidth}{0.4pt}\bigskip

%%--------------------------------------------------------------
%\subsection*{Oktatói tippek -- zárófeladat}
%%--------------------------------------------------------------
%
%\begin{itemize}[noitemsep, topsep=2pt]
%  \item Mindig mutasd be magad előtt a promptot -- az AI néha meglepő
%        választ ad, és ez jó tanítási alkalom.
%  \item Ha gyenge eredmény születik, ne javítsd ki azonnal.
%        Kérdezd meg: ,,Mi hiányzik a promptból?''
%  \item A rossz promptok legalább annyit tanítanak, mint a jók.
%  \item Ha a válasz túl hosszú: ,,Rövidítsd le 3 mondatra.''
%        Ha túl rövid: ,,Fejtsd ki részletesebben, példákkal.''
%  \item \textbf{Hallucináció-teszt:} az A és B dokumentumoknál
%        a konkrét számokat mindig ellenőrizzük az eredetiben.
%        A C dokumentumnál a jegyszámítási képletet érdemes
%        összevetni az eredeti szöveggel.
%  \item \textbf{Összehasonlítás:} az 1.\ feladatban mindhárom
%        változatnál ugyanaz a prompt szerepel -- érdemes megmutatni
%        egymás mellett a három kimenetet. Jól látszik, hogy a
%        dokumentum jellege mennyire befolyásolja az AI válaszát,
%        még azonos prompt esetén is.
%  \item \textbf{Adatvédelem:} személyes adatokat, szerződéseket,
%        bizalmas dokumentumokat ne töltsenek fel publikus Gemini-be.
%\end{itemize}

%=============================================================
\newpage
\section*{Cheat Sheet -- Prompt engineering alapok}
%=============================================================

\begin{tcolorbox}[cheatstyle, title={RTFC keretrendszer}]
\begin{tabularx}{\linewidth}{@{}llX@{}}
\textbf{R} & Role    & \textit{„Te egy pénzügyi elemző vagy\ldots"} \\
\textbf{T} & Task    & \textit{„Foglald össze\ldots", „Elemezd\ldots", „Írj\ldots"} \\
\textbf{F} & Format  & \textit{„max.\ 5 mondat", „Markdown-táblázat", „felsorolás"} \\
\textbf{C} & Context & Adatok, dokumentum, szituáció -- explicit módon elválasztva (pl.\ \texttt{\#\#\#}, \texttt{"""} ) \\
\end{tabularx}
\end{tcolorbox}

\smallskip

\begin{tcolorbox}[cheatstyle, title={Delimiterek}]
\begin{tabular}{@{}ll@{}}
\texttt{\#\#\# Fejléc} & Szakaszok elválasztása (Task, Format, Context) \\[3pt]
\texttt{"""..."""}     & Adatblokk egyértelmű elhatárolása (hosszú szöveg, táblázat, forrás) \\
\end{tabular}

\smallskip\noindent
\textit{Példa: \texttt{\#\#\# Kontextus} majd \texttt{"""BNR inflációs adat\ldots"""}}
\end{tcolorbox}

\smallskip

\begin{tcolorbox}[cheatstyle, title={5 haladó prompttechnika}]
\begin{tabularx}{\linewidth}{@{}lX@{}}
\textbf{Few-shot}         & Adj 2--3 példát a kívánt stílusra, majd kérj új választ hasonló mintában. \\[3pt]
\textbf{Chain of Thought} & \textit{„Gondolkodj lépésről lépésre: 1\ldots{} 2\ldots{} 3\ldots"} \\[3pt]
\textbf{Persona}          & \textit{„Mostantól konzervatív banki hitelbíráló vagy\ldots"} \\[3pt]
\textbf{Negatív utasítás} & \textit{„NE használd: forradalmi, game-changer\ldots"} \\[3pt]
\textbf{Reflexió}         & A prompt végén: \textit{„Nézd át saját válaszodat: vannak-e hibák vagy ellentmondások?"} \\
\end{tabularx}
\end{tcolorbox}

\smallskip

\begin{tcolorbox}[cheatstyle, title={Reflexió -- példa}]
\texttt{\#\#\# Reflexió}\\
\texttt{Nézd át saját válaszodat: vannak-e hibák, ellentmondások?}\\
\texttt{Figyelembe vetted-e a romániai realitásokat?}\\
\texttt{Add meg a javított, végleges verziót.}
\end{tcolorbox}

\smallskip

\begin{tcolorbox}[cheatstyle, title={Hasznos iterációk}]
\begin{tabular}{@{}ll@{}}
\textit{„Legyen formálisabb."}         & \textit{„Rövidítsd 3 mondatra."} \\[2pt]
\textit{„Ne használj szakzsargont."}   & \textit{„Adj hozzá egy táblázatot."} \\[2pt]
\textit{„Fejtsd ki részletesebben."}   & \textit{„Fordítsd román nyelvre."} \\[2pt]
\multicolumn{2}{@{}l}{\textit{„Magyarázd el, mintha 16 éves lennék."}} \\[2pt]
\multicolumn{2}{@{}l}{\textit{„Ha gyenge a válasz: adj több kontextust, ne csak ismételd a kérést."}} \\
\end{tabular}
\end{tcolorbox}

\smallskip

\begin{tcolorbox}[cheatstyle, title={Hallucináció – figyelmeztetés}]
\begin{tabularx}{\linewidth}{@{}lX@{}}
\textbf{Mindig ellenőrizd:} & Számokat, dátumokat, neveket, hivatkozásokat az eredeti dokumentumban. \\[3pt]
\textbf{Dokumentumfeltöltés:} & PDF feltöltése a Gemini rendszerébe $\neq$ a szöveg egyszerű beillesztése --
        a feltöltött fájl tartalmát az AI strukturáltabban kezeli. \\[3pt]
\textbf{Bizonytalanság:} & Ha nem talál információt, az AI hajlamos kiegészíteni vagy feltételezni --
        kérd meg: \textit{„Ha nem szerepel a dokumentumban, írd: nem szerepel."} \\
\end{tabularx}
\end{tcolorbox}

\smallskip

\begin{tcolorbox}[warnstyle]
\textbf{\textcolor{warnborder}{Adatvédelem:}} 
Ne tölts fel bizalmas adatokat, szerződéseket vagy személyes információkat
publikus Gemini rendszerbe.
Használj anonimizált vagy fiktív adatokat.
\end{tcolorbox}
\vfill
\begin{center}
\small\textcolor{medgray}{Sapientia EMTE \textbullet{} Generatív AI és üzleti alkalmazások \textbullet{} 2026}
\end{center}

\end{document} 